\documentclass[../../../include/open-logic-section]{subfiles}

\begin{document}
	
\olfileid[hi]{sth}{z}{nat}
\olsection{अपनी प्राकृतिक संख्याओं का चयन}

%Let us take it for granted that the informal principle \stagesinf{}
%can be justified. It is also, though, worth commenting on the
%particular axiom we extracted from that informal principle. 

\olref[infinity-again]{defnomega} में हमने ``प्राकृतिक संख्याएँ'' पदबंध को
स्पष्टतः परिभाषित किया। इस नियतन को कैसे समझना चाहिए? यह कोई तत्त्वमीमांसात्मक
दावा नहीं, बल्कि कुछ समुच्चयों को प्राकृतिक संख्याएँ \emph{मानने} का निर्णय
मात्र है। ऐसा सोचने के कारणों को हमने \olref[sfr][rel][ref]{sec},
\olref[sfr][arith][ref]{sec} और
\olref[sfr][infinite][dedekindsproof]{sec} में छुआ था। किंतु इन कारणों को और
तीक्ष्ण बनाया जा सकता है।

हमारा अनंतता का अभिगृहीत \citet{VonNeumann1925} का अनुसरण करता है। किंतु
एक दूसरा अभिगृहीत है जिसे हम उसके स्थान पर अपना सकते थे:

\begin{defish}
\emph{ज़ेर्मेलो का \citeyear{Zermelo1908Untersuchungen} का अनंतता का
अभिगृहीत।} ऐसा समुच्चय $A$ है कि $\emptyset \in A$ और
$(\forall x \in A)\{x\} \in A$।
\end{defish}

यदि हमने अपने (वॉन नोयमान-प्रेरित) अनंतता के अभिगृहीत के स्थान पर ज़ेर्मेलो
का अभिगृहीत उपयोग किया होता, तो हमें उतने ही अच्छे ढंग से एक डेडेकिंड-अनंत
समुच्चय और इसलिए एक डेडेकिंड बीजगणित मिलता। ज़ेर्मेलो के दृष्टिकोण में हमारे
बीजगणित का विशिष्ट अवयव फिर $\emptyset$ (हमारा $0$ का प्रतिनिधि) होता, किंतु
अंतःक्षेप $x \mapsto x \cup \{x\}$ के बजाय प्रतिचित्रण
$x \mapsto \{x\}$ से मिलता। इसका सरलतम परिणाम यह है कि ज़ेर्मेलो $2$ को
$\{\{\emptyset\}\}$ मानते हैं, जबकि हम (वॉन नोयमान के साथ) $2$ को
$\{\emptyset, \{\emptyset\}\}$ मानते हैं।

अनंतता के एक अभिगृहीत को दूसरे के बजाय क्यों चुनें? मुख्य व्यावहारिक कारण
यह है कि वॉन नोयमान का दृष्टिकोण परिमितेतर संख्याओं को सँभालने के लिए बहुत
अच्छी तरह ``विस्तारित'' होता है। इसे हम \olref[ordinals][]{chap} से आगे
देखेंगे। किंतु केवल \emph{अंकगणित करने} के दृष्टिकोण से दोनों दृष्टिकोण
समान रूप से अच्छे होते। अतः यदि कोई आपसे कहे कि प्राकृतिक संख्याएँ समुच्चय
\emph{हैं}, तो स्पष्ट प्रश्न है: \emph{वे कौन-से समुच्चय हैं?}

इसी प्रश्न को \citet{Benacerraf1965} ने प्रसिद्ध किया। किंतु इस बात पर बल
देना उचित है कि यह उस परिघटना का केवल सबसे प्रसिद्ध उदाहरण है जिससे हम पहले
ही कई बार मिले हैं। मूल बात यह है। समुच्चय सिद्धांत हमें कई ``सहज'' प्रकार
की सत्ताओं का \emph{अनुकरण} करने का मार्ग देता है: वास्तविक, परिमेय, पूर्ण
और प्राकृतिक संख्याएँ तो हैं ही; साथ ही क्रमित युग्म, फलन और संबंध भी। किंतु
समुच्चय सिद्धांत अनुकरण का कभी कोई \emph{अद्वितीय} चयन नहीं देता। \emph{सदा}
ऐसे विकल्प होते हैं जो---सीधे-सीधे---उतनी ही अच्छी तरह काम करते।

\end{document}
